\documentclass[12pt]{article}

\setlength{\parindent}{0pt}
\setcounter{secnumdepth}{3}
\usepackage[margin=1in]{geometry}

% packages
\usepackage{amsmath, amsfonts, amsthm, amssymb}
\usepackage{graphicx, float}
\usepackage{mathtools}
\usepackage{titlesec}
\usepackage{interval}
\usepackage{hyperref}
\usepackage{siunitx}
\usepackage{titling}
\usepackage{vwcol}
\usepackage{setspace}
\usepackage{empheq}
\usepackage{cancel}
\usepackage{esdiff}
\usepackage{multicol}
\usepackage{mdframed}
\usepackage{esdiff}
\usepackage{tikzsymbols}
\usepackage{multicol}
\usepackage{tikz}
\usepackage{varwidth}
\usepackage{pgfplots}
\usepackage{fancyhdr}
\usepackage{enumitem}
\usepackage{tcolorbox}

\intervalconfig {
	soft open fences
}

% header
\pagestyle{fancy}
\fancyhf{}
\lhead{Ethan Fan}
\rhead{PCH}
\cfoot{\thepage}

% section format
\titleformat{\section}
  {\large \bfseries}
  {}
  {0em}
  {}
\titlespacing*{\section}{0pt}{0.5em}{0.3em}

\titleformat{\subsection}[runin]
  {\bfseries}
  {}
  {0em}
  {}
\titlespacing*{\subsection}{0pt}{0pt}{0em}

\titleformat{\subsubsection}[runin]
  {\itshape}
  {(\roman{subsubsection})}
  {0em}
  {}

% commands
\newcommand{\alignedintertext}[1]{%
  \noalign{%
    \vskip\belowdisplayshortskip
    \vtop{\hsize=\linewidth#1\par
    \expandafter}%
    \expandafter\prevdepth\the\prevdepth
  }%
}

\newcommand*{\paren}[1]{\ensuremath\left(#1\right)}
\newcommand*{\limit}[2][x]{\ensuremath{\displaystyle\lim_{#1\to#2}}}
\newcommand*{\Limit}[3][x]{\ensuremath{\displaystyle\lim_{#1\to#2}\left[#3\right]}}
\newcommand*{\deriv}[1][x]{\ensuremath{\dfrac{\mathrm{d}}{\mathrm{d}#1}}}
\newcommand*{\Deriv}[2][x]{\ensuremath{\dfrac{\mathrm{d}}{\mathrm{d}#1}\left[#2\right]}}
\newcommand{\dydx}{\dfrac{dy}{dx}}
\newcommand{\dydt}{\dfrac{dy}{dt}}
\newcommand{\dxdt}{\dfrac{dx}{dt}}
\newcommand*{\iinteg}[2][x]{\ensuremath{\displaystyle\int #2\;\mathrm{d}#1}}
\newcommand*{\dinteg}[4][x]{\ensuremath{\displaystyle\int_{#2}^{#3}#4\;\mathrm{d}#1}}
\newcommand*{\abs}[1]{\ensuremath{\left|#1\right|}}
\newcommand*{\eps}{\varepsilon}
\newcommand*{\real}{\ensuremath\mathbb{R}}
\newcommand*{\integer}{\ensuremath\mathbb{Z}}
\newcommand*{\posint}{\ensuremath\mathbb{Z^+}}
\newcommand*{\cbrt}[1]{\ensuremath{\sqrt[3]{#1}}}
\newcommand*{\lheqzero}{\ensuremath{\underset{\text{L'H}}{\overset{\left[\frac00\right]}{=}}}}
\newcommand*{\lheqinfty}{\ensuremath{\underset{\text{L'H}}{\overset{\left[\frac{\infty}{\infty}\right]}{=}}}}
\newcommand{\tor}{\text{ or }}
\newcommand{\floor}[1]{\lfloor #1 \rfloor}

\newcommand{\vem}{\vspace{0.5em}}
\newcommand{\example}[1]{\section{Example #1}}
\newcommand{\problem}[1]{\section{Problem #1}}
\newcommand{\subproblem}[1]{\subsection{(#1)} \,}

\DeclareMathOperator{\dne}{DNE}
\DeclareMathOperator{\sgn}{sgn}

\newcommand{\definition}[1]{\begin{tcolorbox}[colback=red!5!white,colframe=red!75!black,parbox=false] #1 \end{tcolorbox}}

\newcommand{\theorem}[2]{\begin{tcolorbox}[title={#1},colback=blue!5!white,colframe=blue!75!black,parbox=false] #2 \end{tcolorbox}}

\allowdisplaybreaks
% \postdisplaypenalty=100000

% document
\begin{document}

\vspace*{0cm}

% opening
\begin{center}
    {\Large \bfseries Problem Set 76} \\[0.5cm]
    {\large Optimization Part 2} \\[0.35cm]
    {\normalsize May 19, 2026}
\end{center}

% problems

\problem{4}
 An observer stands at a point $P$, one unit away from a track. Two runners start at the point in the
figure and run along the track. One runner runs three times as fast as the other. Find the maximum
value of the observer’s angle of sight $\theta$ between the runners. [Hint: Maximize $\tan \theta$.]
\begin{gather*}
    y=\tan \theta\\
    y=\tan (\theta_2-\theta _1)=\frac{3x-x}{1+3x^2}\\
    y=\frac{2x}{1+3x^2},\quad x>0\\
    \dydx=\frac{2(1+3x^2)-12x^2}{(1+3x^2)^2}\\
    \dydx =0 \implies 1=3x^2 \implies \boxed{x=\frac{1}{\sqrt{3}}}
\end{gather*}
As $y'$ changes from positive to negative at $x=\dfrac{1}{\sqrt{3}}$, $y$ has a local maximum at $x=\dfrac{1}{\sqrt{3}}$.
\begin{gather*}
    \tan \theta = \frac{\frac{2}{\sqrt{3}}}{1+3\cdot \frac{1}{3}}=\frac{1}{\sqrt{3}}\implies \boxed{\theta = \frac{\pi}{6}}
\end{gather*}
The maximum value of the observer's angle of sight is $\dfrac{\pi}{6}$.

\problem{5}
 Where should the point P be chosen on the line segment AB so as to maximize the angle $\theta$?
 \begin{gather*}
     \begin{cases}
         x+y=3\\
         x=2\tan \theta _1\\
         y=5\tan \theta_2
     \end{cases}\\
     x=2\tan \theta _1  \implies \theta_1=\arctan \frac{x}{2}\\
     y=5\tan \theta_2 \implies \theta_2 = \arctan \frac{y}{5}\\
     \theta = \theta_1+\theta_2 = \arctan \frac{x}{2}+ \arctan \frac{(3-x)}{5}, \quad 0<x<3\\
     \frac{d\theta}{dx}=\frac{1}{1+\frac{x^2}{4}}\cdot \frac{1}{2}+\frac{1}{1+\frac{x^2-6x+9}{25}}\cdot (-\frac{1}{5})\\ 
     \frac{d\theta}{dx}=0 \implies \frac{1}{1+\frac{x^2}{4}}\cdot \frac{1}{2}=\frac{1}{1+\frac{x^2-6x+9}{25}}\cdot \frac{1}{5}\\
     20+5x^2=50+2x^2-12x+18\\
     x^2+4x-16=0\\
     x=\frac{-4\pm 4\sqrt{5}}{2} \implies \boxed{x=-2+2\sqrt{5}}
 \end{gather*}
 Point $P$ is $5-2\sqrt{5}$ above $A$.

 \problem{6}
Find the maximum area of a rectangle that can be circumscribed about a given rectangle with length $L$
and width $W$. [Hint: Express the area as a function of an angle $\theta$.]
\begin{gather*}
    \begin{cases}
        A=xy\\
        x=L \cos \theta + W\sin \theta\\
        y=L\sin \theta + W\cos \theta
    \end{cases}\\
    A=(L\cos \theta +W\sin \theta )(L\sin \theta + W\cos \theta)\\
    A=(L^2+W^2)\sin \theta \cos \theta + WL (\sin^2\theta + \cos^2\theta)\\
    A=(L^2+W^2)\cdot \frac{\sin 2\theta}{2}+ WL, \quad 0<\theta<\frac{\pi}{2}\\
    A'=(L^2+W^2)\cos 2\theta\\
    A'=0 \implies \theta = \frac{\pi}{4}
\end{gather*}
As $A'$ changes from positive to negative at $\theta = \dfrac{\pi}{4}$, $A$ has a local maximum at $\theta = \dfrac{\pi}{4}$.
\begin{gather*}
    A=(L^2+W^2)\cdot \frac{\sin 2\cdot \frac{\pi}{4}}{2}+ WL=\frac{L^2}{2}+\frac{W^2}{2}+WL\\
    \boxed{A=\frac{1}{2}(W+L)^2}
\end{gather*}
The maximum area of a rectangle is $A=\dfrac{1}{2}(W+L)^2$.

\problem{7}
Ornithologists have determined that some species of birds tend to avoid flights over large bodies of water
during daylight hours. It is believed that more energy is required to fly over water than over land because
air generally rises over land and falls over water during the day. A bird with these tendencies is released
from an island that is 5 km from the nearest point on a straight shoreline, flies to a point on the shoreline,
and then flies along the shoreline to its nesting area. Assume that the bird instinctively chooses a path
that will minimize its energy expenditure. Points and are 13 km apart. \vem

\subproblem{a}
In general, if it takes 1.4 times as much energy to fly over water as it does over land, to what point
$C$ should the bird fly in order to minimize the total energy expended in returning to its nesting area?
\begin{gather*}
    x+y=13\\
    J=\frac{W}{L}\sqrt{x^2+25}+y\\
    J=\frac{W}{L}\sqrt{x^2+25}+13-x,\quad 0<x<13\\
    J'=\frac{W}{L}\frac{x}{\sqrt{x^2+25}}-1\\
    J'=0 \implies L\sqrt{x^2+25}=Wx\\
    L^2x^2+25L^2=W^2x^2\\
    \boxed{x=\sqrt{\frac{25L^2}{W^2-L^2}}}
\end{gather*}
Given $\dfrac{W}{L}=\dfrac{7}{5}$,
\begin{gather*}
    x=\sqrt{\frac{25\cdot 25}{49-25}}=\boxed{\frac{25}{2\sqrt{6}}}
\end{gather*}
The bird should fly to a point $\dfrac{25}{2\sqrt{6}}$ km away from $B$ towards $C$.

\subproblem{b}
Let $W$ and $L$ denote the energy (in joules) per kilometer flown over water and land, respectively.
What would a large value of the ratio $W/L$ mean in terms of the bird’s flight? What would a small value
mean? Determine the ratio $W/L$ corresponding to the minimum expenditure of energy. \vem

As attained from part $\textbf{(a)}$:
\[
x=\sqrt{\frac{25L^2}{W^2-L^2}}=\sqrt{\frac{25}{(\frac{W}{L})^2-1}}
\]
A larger value of the ratio $W/L$ indicates that the bird first flies to a point on the shoreline closer to $B$, while a smaller value indicates that the bird first flies to a point on the shoreline closer to $C$. \vem

Regarding the ratio of $W/L$ with the minimum expenditure of energy:
\begin{gather*}
    J'=0 \implies \frac{W}{L}\cdot \frac{x}{\sqrt{x^2+25}}-1=0\\
    \boxed{\frac{W}{L}=\frac{\sqrt{x^2+25}}{x}}
\end{gather*}

\subproblem{c}
What should the value of $W/L$ be in order for the bird to fly directly to its nesting area $D$? What
should the value of $W/L$ be for the bird to fly to $B$ and then along the shore to $D$? \vem

Plugging in the equation from part $\textbf{(b)}$ gives:
\begin{gather*}
    x=13 \implies \boxed{\frac{W}{L}=\frac{\sqrt{194}}{13}}\\
    x=0 \implies \boxed{\frac{W}{L} \to \infty}
\end{gather*}

\subproblem{d}
If the ornithologists observe that birds of a certain species reach the shore at a point 4 km from $B$,
how many times more energy does it take a bird to fly over water than over land?
\begin{gather*}
    x=4 \implies \boxed{\frac{W}{L}=\frac{\sqrt{41}}{4}}
\end{gather*}



\end{document}

